% ============================================
% DOC-22
% Extraído de main.pandoc.tex
% ============================================
\fragtitle{doc}{D}{DOC}{22}{Memorando Interno nº~47/2026}
\markboth{DOC-22 --- Memorando Interno nº~47/2026}{}

\noindent \textit{Documento extraído do arquivo pericial.}

\vspace{0.5em}\noindent\rule{\textwidth}{0.5pt}\vspace{0.5em}

\noindent {\footnotesize\texttt{ARQUIVO P\'UBLICO $\cdot$ PROVA DOCUMENTAL $\cdot$ PROCESSO 1.260}}\\[-0.2em]

\noindent \textbf{Arquivo Público Mineiro --- Diretoria de Documentação}\\
\noindent \textbf{Para:} Setor de Acervos Especiais. \textbf{CC:} Superintendência Administrativa.\\
\noindent \textbf{Assunto:} Encerramento de análise do objeto catalogado sob registro 1.260.\\
\noindent\textbf{Data:} 1º de agosto de 2026.

\noindent Prezados,

\noindent 1. Considerando que o objeto caderno atribuído ao registro 1.260 (Hospital Colônia de Barbacena) foi submetido a análise externa pela psiquiatra forense Dra. Lúcia Mendes (CRM-MG 28.391), e que a referida profissional solicitou afastamento clínico antes da conclusão do parecer;

\noindent 2. Considerando que o laudo de re-diagnóstico (DOC-08) já foi anexado ao prontuário original e que não há, neste momento, determinação judicial, administrativa ou acadêmica que justifique a manutenção do objeto em circulação;

\noindent 3. Considerando as diretrizes de segurança documental da instituição e o princípio de estabilidade do acervo;

\noindent \textbf{determino:}

\noindent a) O encerramento imediato da consulta pública ao objeto 1.260;

\noindent b) A guarda do objeto em cofre definitivo, sem previsão de novo acesso;

\noindent c) A notificação de todas as partes envolvidas de que o caso está arquivado e não serão admitidas novas solicitações de manuseio.

\noindent \textbf{Fundamentação:} o objeto não constitui prova material de crime, não integra ação judicial em curso e não se qualifica como documento de valor histórico excepcional que justifique risco ao corpo técnico. A permanência de sua circulação atenta contra a segurança documental e a integridade psicológica dos profissionais envolvidos.

\noindent \textbf{Recomendação adicional:} sugiro que a diretoria avalie a conveniência de não registrar este memorando no sistema informatizado de gestão documental. O registro informatizado torna o documento localizável. E um documento localizável pode ser solicitado novamente.

\noindent \textit{Ronaldo Torres}\\
\noindent \textit{Arquivo Público Mineiro --- Diretor de Documentação}
\noindent \textit{Anotação manuscrita na margem superior, tinta azul, letra cursiva:}
\noindent \textquotedblleft{}\textit{Ronaldo --- não foi possível. O objeto foi retirado do cofre na manhã seguinte ao memorando. O sistema de rastreamento registra a abertura do cofre às 3h14. Ninguém foi identificado nas câmeras. A porta estava lacrada. O lacre estava íntegro. Assunto arquivado conforme solicitado.\textquotedblright{}}

\noindent \textquotedblleft{}\textit{--- Setor de Acervos Especiais, 02/08/2026.\textquotedblright{}}

\vspace{0.5em}
\noindent\textit{\textbf{↪ Rotas:} \rota{DOC-01} \rotasep \rota{LUC-02} \rotasep \rota{DOC-23}}